% ---------------------------------------------------------------------
%  附录 A 公式汇总
% ---------------------------------------------------------------------
\chapter{公式汇总}
\label{app:formulas}

本附录把正文中反复使用的公式集中列出，便于查阅。

\section{轮廓与分量}
\begin{align}
    G(t,t') &= -\ii\,\langle \mathcal{T}_{\Contour}
        c(t)c^{\dagger}(t')\rangle,\\
    \Gret(t,t') &= -\ii\theta(t-t')
        \langle[c(t),c^{\dagger}(t')]_{\mp}\rangle,\\
    \Gless(t,t') &= \mp\ii\,\langle
        c^{\dagger}(t')c(t)\rangle,\\
    \Glmix(t,\tau) &= \mp\ii\,\langle
        c^{\dagger}(\tau)c(t)\rangle,\\
    \Gmat(\tau) &= -\langle
        \mathcal{T}_\tau c(\tau)c^{\dagger}(0)\rangle .
\end{align}

其中 $\mp$ 与 $[\cdot,\cdot]_{\mp}$ 的取法为：玻色型取上面的符号
（$-$，对易子），费米型取下面的符号（$+$，反对易子）。

\section{共轭与重构}
\begin{align}
    \Gadv(t,t') &= \big[\Gret(t',t)\big]^{\dagger},\\
    \Gless(t,t') &= -\big[\Gless(t',t)\big]^{\dagger},\\
    \Grmix(\tau,t) &= \mp\big[\Glmix(t,\bet-\tau)\big]^{\dagger},\\
    \tilde{G}^{R}_{j,m} &=
        -\big([G^{\ddagger}]^{R}_{m,j}\big)^{\dagger},
        \qquad j<m .
\end{align}

\section{离散卷积}
\begin{align}
    C^{R}(n,m) &=
        h\sum^{n}_{j=m}w^{(k)}_{n-m,\,j-m}
        A^{R}_{n,j}f_j B^{R}_{j,m},
        && n>k,\ n-m>k,\\
    C^{R}(n,m) &=
        h\sum^{k}_{j=0}w^{(k)}_{n-m,\,j}
        A^{R}_{n,n-j}f_{n-j}\tilde{B}^{R}_{n-j,m},
        && n>k,\ n-m\le k,\\
    C^{R}(n,m) &=
        h\sum^{k}_{j=0}I^{(k)}_{m,n;j}
        \tilde{A}^{R}_{n,j}f_j\tilde{B}^{R}_{j,m},
        && n\le k .
\end{align}

\section{权重}
Gregory 权重 $\Omm^{(k)}_{nj}$ 的五种情形见 \cref{eq:omega-cases}。
其余三张表由多项式插值导出：
\begin{align}
    a^{(k)}_j &= -D^{(k)}_{0j},\\
    D^{(k)}_{ml} &= \sum^{k}_{a=1}P^{(k)}_{al}\,a\,m^{a-1},\\
    P^{(k)}_{al} &= (M^{-1})_{al},\qquad M_{ja}=j^{a}.
\end{align}

\section{求解器}
VIE2 的右作用方程（推迟分量，逐列推进）：
\begin{equation}
    \Gret(n,j)\big[1+h\,\omega_0\,\bm{K}^{R}(j,j)\big]
    =
    \bm{Q}^{R}(n,j)
    -h\sum_{m>j}w_m\,\Gret(n,m)\,\bm{K}^{R}(m,j),
\end{equation}
左混合分量的整行形式：
\begin{equation}
    \big[\bm{1}+h\,\Omm_{n0}\,\bm{F}^{R}(n,n)\big]\Glmix(n,\tau)
    =
    \bm{Q}^{\rceil}(n,\tau)-[\bm{F}*G]^{\rceil}_{\mathrm{known}}(n,\tau).
\end{equation}
Dyson 的左混合分量（含 BDF 时间导数）：
\begin{equation}
    \Big[\frac{\ii}{h}b_0+\mu-\bm{H}_n-h\,\Omm_{n0}\,\Sigret(n,n)\Big]
    \Glmix(n,\tau)
    =
    [\Sig*G]^{\rceil}_{\text{known}}(n,\tau)
    -\frac{\ii}{h}\sum^{k+1}_{p=1}b_p\,\Glmix(n-p,\tau).
\end{equation}

\section{Matsubara}
\begin{align}
    \omega_m &= \frac{\pi}{\bet}(2m+\delta_\xi),\\
    \Gmat(\ii\omega_m) &=
    \big[\ii\omega_m+\mu-\bm{H}
      -\Sig^{\mathrm{M}}(\ii\omega_m)\big]^{-1},\\
    c_0(\tau) &=
    \begin{cases}
        -a_0/2, & \text{费米型},\\
        a_0(\tau/\bet-1/2), & \text{玻色型}.
    \end{cases}
\end{align}

\section{稳态与移动窗}
\begin{align}
    \Gret(\omega) &= \big[\omega+\mu-\bm{H}
        -\Sigret(\omega)-\Sigret_b(\omega)\big]^{-1},\\
    \Gless(\omega) &= \Gret(\omega)
        \big[\Sigless(\omega)+\Sigless_b(\omega)\big]
        \Gret(\omega)^{\dagger},\\
    \Delta\omega &= \frac{\pi}{(N_{ft}/2)h},\\
    E_{\mathrm{kin}}(t) &=
        \mathrm{Im}\,\Tr\Big[\bm{h}
        \sum^{t_c}_{l=0}w_l
        \big(\Gret(t,t-l)^{\dagger}\Gless(t,t-l)
        -\Gless(t,t-l)^{\dagger}\Gret(t,t-l)\big)\Big].
\end{align}
移动窗的块 $(\delti,\deltj)$ 保存
$G(t_0-\delti h,\,t_0-\delti h-\deltj h)$。
